\begin{onehalfspace}

Il capitolo traduce la motivazione introduttiva in un contratto sperimentale: definisce ciò che viene verificato, il perimetro entro cui valgono i risultati e le metriche usate. La motivazione crittografica, i requisiti funzionali e la specifica completa di input e output sono riportati nella Sezione~\ref{app:obiettivi} (Sez.~\ref{sec:motivazione}, \ref{sec:requisiti} e~\ref{sec:input_output}).

\section{Obiettivo e Perimetro del Lavoro}

L'obiettivo generale è stabilire quali strategie rendano più affidabile l'esecuzione di Shor in presenza di rumore, distinguendo tre livelli che non sono direttamente intercambiabili: recupero dei candidati dall'istogramma finale, correzione basata sulle sindromi e analisi di sensibilità a un errore logico residuo.

Il lavoro combina tre attività essenziali:
\begin{enumerate}
    \item richiamare il period finding, la \ac{QFT}, la stima di fase e i principali canali di rumore necessari a interpretare gli esperimenti;
    \item implementare e validare Shor, quindi confrontare su simulatori rumorosi la baseline TOP-1, la ricerca TOP-4 e la variante con classificatore, usando gli stessi istogrammi;
    \item costruire una progressione QEC --- ripetizione, Steane $[[7,1,3]]$ e surface code --- e misurare quando la codifica o un decoder appreso riducono l'errore logico. Una successiva analisi applica a Shor un proxy Pauli per gate compilato, senza identificarlo con il tasso di un memory experiment.
\end{enumerate}

Le campagne rumorose complete riguardano l'istanza $N=15$. Le istanze $N=21$ e $N=35$ sono impiegate soltanto per la validazione ideale dell'aritmetica e per l'audit delle risorse: non vengono usate per inferire la scalabilità sotto rumore né per rappresentare una QPU reale.

\section{Domande di Ricerca e Ipotesi}
\label{sec:ipotesi}

La domanda trasversale della prima fase è se il classificatore riduca il numero di iterazioni oltre quanto già ottenibile esaminando più candidati dello stesso istogramma. L'ipotesi principale richiede pertanto un vantaggio di M2 sia rispetto alla baseline M1 sia rispetto all'ablazione TOP-4; F1 e AUC del classificatore sono controlli predittivi secondari, non prove autonome di utilità operativa. Un eventuale vantaggio deve inoltre essere osservabile nei due preset $N=15$, mentre la generalizzazione al crescere di $N$ resta fuori dal perimetro rumoroso.

La progressione sperimentale risponde poi alle seguenti domande:
\begin{description}
    \item[\textbf{DR1.}] Come varia l'output di Shor al crescere dei diversi parametri di rumore, e quali conclusioni sono sostenute dagli sweep controllati?
    \item[\textbf{DR2.}] Come si estrae una sindrome con un codice a ripetizione senza misurare direttamente lo stato logico?
    \item[\textbf{DR3.}] Quali vantaggi e limiti offre Steane $[[7,1,3]]$ nella correzione di un errore Pauli arbitrario?
    \item[\textbf{DR4.}] Il surface code mostra un regime di soglia nel quale aumentare la distanza riduce il logical error rate?
    \item[\textbf{DR5.}] Come varia il successo di Shor sotto un proxy Pauli uniforme per gate e quale beneficio producono ripetizioni indipendenti?
    \item[\textbf{DR6.}] In quale regime un decoder appreso dalla sindrome può migliorare il decoder analitico di riferimento?
\end{description}

Le ipotesi quantitative associate sono: soppressione quadratica dell'errore per i codici di distanza $d=3$ nel regime perturbativo; esistenza di un incrocio delle curve $p_L(p)$ del surface code; andamento non crescente del successo di Shor al crescere del proxy per gate. Per la decodifica appresa il confronto è deliberatamente esplorativo: il modello deve essere valutato sugli stessi campioni e contro \ac{MWPM}, senza presupporne la superiorità.

\section{Contributo e Disegno Sperimentale}
\label{sec:contributo_originale}

Il contributo metodologico è costituito da tre scelte coordinate. Primo, il confronto appaiato M1--TOP-4--M2 separa il valore della ricerca multi-candidato da quello del classificatore. Secondo, la progressione dai codici elementari al surface code usa metriche esplicite e dichiara il confine tra tasso per memoria, tasso per ciclo e proxy per gate. Terzo, il modello appreso viene valutato in due punti diversi della pipeline --- output e sindrome --- così che la sua utilità dipenda dall'informazione effettivamente disponibile, non dalla sola qualità predittiva.

Il protocollo mantiene distinti risultati, validazioni e assunzioni. Circuiti e aritmetica sono prima verificati idealmente; i metodi decisionali ricevono gli stessi campioni; ogni stima riporta seed, numerosità e incertezza; i preset uniformi sono trattati come scenari illustrativi e non come calibrazioni hardware. Il dettaglio implementativo e statistico è nel Capitolo~\ref{chap:metodologia}, mentre gli esperimenti estesi sono raccolti nelle Appendici~\ref{app:fase1} e~\ref{app:fase2}.

\section{Campagne Rumorose e Validazioni Ideali}
\label{sec:use_case}

Il disegno distingue due use case:
\begin{itemize}
    \item \textbf{campagne rumorose N15}: stesso $N=15$, base $a=7$, circuito e budget di shot, con un preset di riferimento e uno di stress;
    \item \textbf{validazioni ideali N21/N35}: verifica dell'aritmetica reversibile, del ripristino delle ancille e della distribuzione QPE, senza simulazione rumorosa completa.
\end{itemize}

\begin{table}[!htbp]
\centering
\small
\begin{tabular}{@{}lcc@{}}
\toprule
\textbf{Parametro} & \textbf{Riferimento} & \textbf{Stress} \\
\midrule
$N$, $a$, periodo atteso $r$ & $15,7,4$ & $15,7,4$ \\
$n_{\mathrm{count}}$ & 8 & 8 \\
$M_{\mathrm{shots}}$ & 1024 & 1024 \\
$\epsilon_{1q}=\lambda_{1q}$ & $1\times10^{-3}$ & $5\times10^{-3}$ \\
$\epsilon_{2q}=\lambda_{2q}$ & $1\times10^{-2}$ & $5\times10^{-2}$ \\
$T_1$ (ns) & 100\,000 & 50\,000 \\
$T_2$ (ns) & 80\,000 & 30\,000 \\
$p_{\mathrm{ro}}$ & $2\times10^{-2}$ & $5\times10^{-2}$ \\
\bottomrule
\end{tabular}
\caption{Preset uniformi delle campagne rumorose N15; non rappresentano la calibrazione di una QPU.}
\label{tab:use_case_params}
\end{table}

I parametri $\lambda_{1q}$ e $\lambda_{2q}$ sono quelli passati al canale depolarizzante di Aer; il readout è modellato mediante una matrice di confusione simmetrica. Il modello è uniforme e non comprende topologia, routing, crosstalk o deriva temporale.

\subsection{Razionale per la Scelta delle Istanze}
\label{subsec:razionale_istanze}

Per $N=15$, la base $a=7$ ha periodo atteso $r=4$ e consente un circuito abbastanza contenuto da sostenere confronti appaiati e sweep. Per $N=21$, $a=2$, e $N=35$, $a=6$, l'aritmetica di Beauregard è invece validata idealmente sui vettori di base e mediante QPE. La separazione evita di trasformare un limite computazionale del simulatore in un'affermazione sul successo dell'algoritmo.

\section{Metriche di Valutazione}
\label{sec:metriche}

\subsection{Post-processing dell'Output}

M1 usa il solo candidato più frequente; $M_{\mathrm{TOP4}}$ verifica i quattro candidati dominanti; M2 applica il classificatore prima della stessa ricerca TOP-4. Per ciascun metodo $j$ si riportano il numero medio di iterazioni $\bar M_j$, la dispersione e la probabilità di successo entro il budget preregistrato $M_{\max}=50$. Il confronto descrittivo tra baseline e metodo con classificatore usa
\begin{equation}
    \rho=\frac{\bar M_1}{\bar M_2},
    \label{eq:ipotesi_rho}
\end{equation}
mentre l'ablazione richiede il confronto separato con $M_{\mathrm{TOP4}}$. Un valore $\rho>1$ indica soltanto una media inferiore per M2: l'inferenza è condotta sulle differenze appaiate complete. I fallimenti sono conservati mediante una sentinella oltre il budget; test, gestione degli zeri e tie-break deterministico sono specificati nel Capitolo~\ref{chap:metodologia}. F1 e AUC descrivono il classificatore, ma non sostituiscono la metrica operativa.

\subsection{Correzione d'Errore}
\label{subsec:metriche_qec}

La metrica centrale è il \textbf{logical error rate} $p_L$, stimato con intervallo di confidenza in funzione dell'errore fisico $p$ e della distanza $d$. La pendenza $\gamma$ in $p_L\propto p^\gamma$ verifica la soppressione a basso errore; il \emph{break-even} è il regime $p_L<p$; la soglia $p_{th}$ è individuata dall'incrocio delle curve a distanze diverse. Per l'analisi di Shor si misurano inoltre la probabilità per singola esecuzione e quella cumulativa $P(k)=1-(1-P_s)^k$.

Il simbolo $p_L$ compare in due contesti che restano distinti: nei benchmark QEC è un tasso riferito allo specifico circuito di memoria o correzione; nel proxy di Shor è la probabilità di un Pauli non identità dopo ciascun gate compilato incluso nel modello. Una conversione tra i due richiederebbe una compilazione fault-tolerant per operazione e non viene assunta.

\section{Approccio Metodologico}

La trattazione procede dalla teoria alla verifica ideale, quindi alle campagne appaiate e alla QEC. Architettura software, modelli di rumore, protocolli statistici e implementazione dei metodi sono descritti nel Capitolo~\ref{chap:metodologia}; questo capitolo ne fissa soltanto obiettivi, perimetro e criteri di valutazione.

\end{onehalfspace}
